\chapter{Volumes under bounded mass}
\label{chap:mass-bound-volumes}

Bounded mass makes singularity volumes finite, so a deficit between them can control stretching and Lelong numbers. At zero deficit, stretching identifies relative full mass with an envelope condition. Domination then shows that ordinary envelopes preserve lower volume, and upper volume when the lower volume is positive. These comparisons give continuity for decreasing model potentials, exact trace realization under the mass assumptions, and a volume criterion for equality of ordinary and multiplier-ideal envelopes.

Assume throughout this chapter that $X$ has bounded mass, and fix a big smooth real $(1,1)$-form $\theta$. We use the volumes of \cref{def:singularity-volumes}.

\section{Comparison of singularity volumes}
\label{sec:bounded-mass-volume-comparison}

Choose a Hermitian form $\Omega\geq\theta$. The estimates in \cref{prop:bmp-finite,thm:arbitrary-extended-ma} bound the mass of every $\theta$-plurisubharmonic potential by $v_+(\Omega)$. Hence, for $u\in\PSH(X,\theta)$,
\[
 0\leq\underline{\operatorname{vol}}_\theta(u)
 \leq\int_X\mathrm{MA}_\theta(u)
 \leq\overline{\operatorname{vol}}_\theta(u)<+\infty.
\]
If $\theta$ is closed, both volumes equal $\int_X\mathrm{MA}_\theta(u)$, by invariance of non-pluripolar mass under equivalent singularities in the closed theory \cite{SGPT}. For a non-closed background this equality can fail even at $u=0$, as \cref{ex:torus-mass} shows.

Volume increases when the background increases or the singularities decrease. The rooftop and maximum constructions preserve the required ordinary types; locality confines their errors to sets with vanishing mass.

\Needspace{12\baselineskip}
\begin{theorem}\label{thm:volume-monotonicity}
Let $\theta_1,\theta_2$ be smooth real $(1,1)$-forms on $X$. Suppose $\theta_1$ is big, $\theta_1\leq\theta_2$, and $u_i\in\PSH(X,\theta_i)$ satisfy $u_1\preceq u_2$. Then
\[
 \underline{\operatorname{vol}}_{\theta_1}(u_1)
 \leq\underline{\operatorname{vol}}_{\theta_2}(u_2),\qquad
 \overline{\operatorname{vol}}_{\theta_1}(u_1)
 \leq\overline{\operatorname{vol}}_{\theta_2}(u_2).
\]
\end{theorem}
\begin{proof}
Shift $u_1$ so that $u_1\leq u_2$. Choose $v_2\simeq u_2$ in $\PSH(X,\theta_2)$ and, for real $t>0$, set
\[
 p_t=P_{\theta_1}(u_1+t,v_2).
\]
For $t>0$, a sufficiently low constant translate of $u_1$ is a competitor; also $p_t\leq u_1+t$. Thus $p_t\simeq u_1$. The contact theorem, applied in each background on the corresponding contact part, and positivity of the mixed products give
\[
 \mathrm{MA}_{\theta_1}(p_t)
 \leq\mathbf1_{\{u_1\leq v_2-t\}}\mathrm{MA}_{\theta_1}(u_1)
      +\mathrm{MA}_{\theta_2}(v_2).
\]
Indeed $p_t$ is also $\theta_2$-plurisubharmonic and is below $v_2$, so the contact inequality in that background applies. The first measure on the right has mass tending to zero: its indicators tend to zero off the pole sets, and the measure is finite and non-pluripolar. Integrate, let $t\to+\infty$, and then take the infimum over $v_2$.

For the upper volume, choose any $v_1\simeq u_1$ in $\PSH(X,\theta_1)$ and put $w_t=\max(v_1,u_2-t)$. It is $\theta_2$-plurisubharmonic and $w_t\simeq u_2$. On $D_t=\{v_1>u_2-t\}$, locality and mixed-product positivity imply
\[
 \mathbf1_{D_t}\mathrm{MA}_{\theta_1}(v_1)
 \leq\mathbf1_{D_t}\mathrm{MA}_{\theta_2}(w_t).
\]
The sets $D_t$ exhaust the complement of a pluripolar set. Hence integration and passage to the limit give
$\int_X\mathrm{MA}_{\theta_1}(v_1)\leq
\overline{\operatorname{vol}}_{\theta_2}(u_2)$.
Take the supremum over $v_1$.
\end{proof}

Adding a positive current also increases both volumes. For the lower volume, subtract its potential inside an envelope to recover a competitor in the original background.

\begin{proposition}\label{prop:volume-addition}
Let $u\in\PSH(X,\theta)$ and $v\in\PSH(X,\eta)$ for a smooth real $(1,1)$-form $\eta$ on $X$. Then, for either the lower or the upper volume,
\[
 \operatorname{vol}_\theta(u)
 \leq\operatorname{vol}_{\theta+\eta}(u+v).
\]
\end{proposition}
\begin{proof}
The form $\theta+\eta$ is big, by adding the positive current $\eta+\ddc v$ to a strictly positive current in the class of $\theta$. For the upper volume, add $v$ to each competitor for $u$. Positivity of mixed products shows that the resulting mass is at least the original mass.

For the lower volume, choose $w\simeq u+v$ in $\PSH(X,\theta+\eta)$, and let $p=P_\theta(w-v)$. Since $w-v=u+O(1)$ outside the pole sets, the envelope is genuine and $p\simeq u$. The contact theorem says that its measure is carried by $D=\{p+v=w\}$.

\Needspace{8\baselineskip}
Furthermore $p+v\leq w$ everywhere by \cref{lem:qpsh-almost-everywhere-order}. Thus
\[
 \mathrm{MA}_\theta(p)
 \leq\mathbf1_D\mathrm{MA}_{\theta+\eta}(p+v)
 \leq\mathbf1_D\mathrm{MA}_{\theta+\eta}(w).
\]
Integrating and taking the infimum over $w$ proves the assertion.
\end{proof}

\section{Volume deficits and relative full mass}

The upper-volume deficit measures how far one potential can be stretched away from another. The positive lower volume of the less singular potential prevents the stretching envelopes from collapsing. We use $Q_b$ and its competitor family from \cref{def:relative-envelope-class}.

\begin{lemma}\label{lem:volume-deficit-stretch}
Let $u,v\in\PSH(X,\theta)$ satisfy $u\preceq v$, and put
\[
 \delta=\underline{\operatorname{vol}}_\theta(v)>0,
 \qquad D=\overline{\operatorname{vol}}_\theta(v)
             -\overline{\operatorname{vol}}_\theta(u).
\]
If $b>1$ and $b^nD<\delta$, then $Q_b(u,v)$ is genuine.
\end{lemma}
\begin{proof}
Adding constants to $u$ and $v$ only translates the stretching envelope. Normalize $v\leq0$, and write $V=\overline{\operatorname{vol}}_\theta(v)$. Since $b^nD<\delta$, choose $w\simeq u$ in $\PSH(X,\theta)$ with
\begin{equation}\label{eq:relative-volume-deficit}
 V-\int_X\mathrm{MA}_\theta(w)<b^{-n}\delta.
\end{equation}
Translate $w$ so that $w\leq v$. For positive integers $j$, put
\[
 w_j=\max(w,v-j),\qquad h_j=Q_b(w_j,v).
\]
The competitor $v-bj$ and the defining inequality give $v-bj\leq h_j\leq v$. Thus $w_j,h_j\simeq v$, and $h_j$ decreases. If its limit were $-\infty$, \cref{thm:normalized-compactness} would give $\sup_Xh_j\to-\infty$.

By \cref{thm:big-general-envelope-contact}, the measure of $h_j$ is carried by $C_j=\{b^{-1}h_j+(1-b^{-1})v=w_j\}$. Ordered contact and positivity of mixed products give
\[
 \mathrm{MA}_\theta(h_j)
       \leq b^n\mathbf1_{C_j}\mathrm{MA}_\theta(w_j).
\]
Fix a positive integer $k$. For large $j>k$, collapse would imply $h_j\leq-bk$ everywhere. Since $v\leq0$, contact then implies $w_j\leq-k$, hence $w\leq-k$. Locality on $\{w>-k\}\subset\{w>v-j\}$ gives
\[
 \begin{split}
 \delta\leq\int_X\mathrm{MA}_\theta(h_j)
 &\leq b^n\int_{\{w\leq-k\}}\mathrm{MA}_\theta(w_j)\\
 &\leq b^n\left(V-\int_{\{w>-k\}}\mathrm{MA}_\theta(w)\right).
 \end{split}
\]
Here the first inequality uses $h_j\simeq v$, and the last uses $w_j\simeq v$. Letting $k\to\infty$ contradicts \eqref{eq:relative-volume-deficit}, since the finite non-pluripolar measure $\mathrm{MA}_\theta(w)$ gives zero mass to the pole set. The limit $h$ is therefore genuine and satisfies $b^{-1}h+(1-b^{-1})v\leq w$. If $w\leq u+C$, then $h-bC\in\mathcal A_b(u,v)$, proving the assertion.
\end{proof}

\begin{proposition}\label{prop:volume-deficit-lelong}
Fix a smooth compact modification $\pi:Z\to X$. There is a finite constant $C_\pi$, depending only on $\pi$ and $\theta$, such that, for $u,v,\delta,D$ as in \cref{lem:volume-deficit-stretch},
\[
 0\leq\nu(u\circ\pi,z)-\nu(v\circ\pi,z)
 \leq C_\pi\min\left\{1,\left(\frac D\delta\right)^{1/n}\right\}
 \qquad(z\in Z).
\]
The same estimate holds for generic Lelong numbers along every prime divisor on $Z$.
\end{proposition}
\begin{proof}
Apply \cref{lem:gauduchon-uniform-lelong} to the smooth background $\pi^*\theta$ on $Z$, and choose $C_\pi$ bounding all its Lelong numbers. Order gives the lower bound and the upper bound $C_\pi$, which proves the assertion when $D\geq\delta$.

If $D<\delta$, every $b>1$ with $b^nD<\delta$ admits $h\in\mathcal A_b(u,v)$ by the lemma. Pulling back the defining convex inequality and taking Lelong numbers gives
\[
 \nu(u\circ\pi,z)-\nu(v\circ\pi,z)
 \leq b^{-1}\bigl(\nu(h\circ\pi,z)-\nu(v\circ\pi,z)\bigr)
 \leq b^{-1}C_\pi.
\]
For $D>0$, let $b\uparrow(\delta/D)^{1/n}$; for $D=0$, let $b\to\infty$. Taking generic values proves the divisorial assertion.
\end{proof}

At zero deficit, every finite stretch is possible. Fix a reference $\varphi\in\PSH(X,\theta)$.

\begin{definition}\label{def:relative-volume-class}
Define
\[
 \mathcal E_{\mathrm{vol}}(X,\theta,\varphi)
 =\{u\in\PSH(X,\theta):u\preceq\varphi,
 \ \overline{\operatorname{vol}}_\theta(u)
   =\overline{\operatorname{vol}}_\theta(\varphi)\}.
\]
\end{definition}

Convex comparison turns the envelope condition back into equality of volumes.

\begin{theorem}\label{thm:relative-equivalence}
Suppose $\underline{\operatorname{vol}}_\theta(\varphi)>0$. Then
\[
 \mathcal E_{\mathrm{env}}(X,\theta,\varphi)
 =\mathcal E_{\mathrm{vol}}(X,\theta,\varphi).
\]
Every member also satisfies
$\underline{\operatorname{vol}}_\theta(u)
=\underline{\operatorname{vol}}_\theta(\varphi)$.
\end{theorem}
\begin{proof}
First assume the envelope condition. For each $a\in(0,1)$ there is $h\in\PSH(X,\theta)$ such that $a\varphi+(1-a)h\leq u$. By \cref{thm:volume-monotonicity,prop:volume-addition} and homogeneity, for either volume,
\[
 a^n\operatorname{vol}_\theta(\varphi)
 \leq\operatorname{vol}_\theta(a\varphi+(1-a)h)
 \leq\operatorname{vol}_\theta(u)
 \leq\operatorname{vol}_\theta(\varphi).
\]
Let $a\to1$. This proves equality of both volumes in this direction, even without positive lower volume.

Conversely, equality of the upper volumes makes $D=0$ in \cref{lem:volume-deficit-stretch}, with $v=\varphi$. Every $Q_b(u,\varphi)$ is genuine, which is the envelope condition.
\end{proof}

\section{Domination and ordinary envelopes}

Positive lower volume prevents comparison envelopes from collapsing. This gives a domination principle and identifies the limit of relative obstacle envelopes, which in turn compares a potential's volumes with those of its ordinary envelope. Write $\mu_w=\mathrm{MA}_\theta(w)$.

\begin{theorem}\label{thm:big-relative-domination}
Let $\varphi\in\PSH(X,\theta)$ have positive lower volume, let $u\in\mathcal E_{\mathrm{env}}(X,\theta,\varphi)$, $v\in\PSH(X,\theta)$ with $v\preceq\varphi$, and $0\leq c<1$. If
\[
 \mathbf1_{\{u<v\}}\mu_u\leq c\,\mathbf1_{\{u<v\}}\mu_v,
\]
then $u\geq v$.
\end{theorem}
\begin{proof}
Put $\ell=\underline{\operatorname{vol}}_\theta(\varphi)>0$ and $\mathcal E_\varphi=\mathcal E_{\mathrm{env}}(X,\theta,\varphi)$. Replace $v$ by $\max(u,v)$. Locality preserves the hypothesis and upward stability puts the new $v$ in $\mathcal E_\varphi$. Translate $u,v$ by the same constant so that $v\leq\varphi$, and then translate all three potentials so that $\varphi\leq0$. Set
\[
 q_b=Q_b(u,v),\qquad
 D_b=\{b^{-1}q_b+(1-b^{-1})v=u\}.
\]
\Cref{lem:relative-stretch-closure} gives $q_b\in\mathcal E_\varphi$ and $q_b\leq u$. Contact support and positivity of the mixed products give
\[
 b^{-n}\mu_{q_b}+(1-b^{-1})^n\mathbf1_{D_b}\mu_v
 \leq\mathbf1_{D_b}\mu_u.
\]
Here $\mu_{q_b}$ is supported on $D_b$. Choose $b$ so large that $(1-b^{-1})^n>c$. Subtracting finite measures on $D_b\cap\{u<v\}$ shows that $\mu_{q_b}$ vanishes there. Thus it is supported on $D_b\cap\{u=v\}\subset\{q_b=u\}$. Since $q_b\leq u$, ordered contact with coefficient one and \cref{thm:relative-equivalence} now give
\begin{equation}\label{eq:big-domination-fixed-tail}
 \ell\leq\mu_{q_b}(X)\leq\mu_u(\{q_b=u\}).
\end{equation}

If $u<v$ anywhere, there is $\delta>0$ such that $\{u<v-\delta\}$ has positive smooth volume. Indeed, the contrary would give $u\geq v$ almost everywhere and then everywhere by \cref{lem:qpsh-almost-everywhere-order}. On this set
$q_b\leq u-(b-1)\delta$. The family $q_b$ decreases in $b$ and is bounded above by $u$. A genuine limit could not be $-\infty$ on a set of positive smooth volume. Hence its limit is identically $-\infty$, and \cref{thm:normalized-compactness} gives $\sup_Xq_b\to-\infty$. The sets $\{q_b=u\}$ then have $\mu_u$-mass tending to zero by dominated convergence off the pole set of $u$. This contradicts \eqref{eq:big-domination-fixed-tail}.
\end{proof}

\begin{corollary}\label{cor:big-relative-orders}
For $u,v$ as in \cref{thm:big-relative-domination} and $\lambda>0$,
\[
 \mathrm{e}^{-\lambda u}\mu_u\leq\mathrm{e}^{-\lambda v}\mu_v
 \quad\Longrightarrow\quad u\geq v.
\]
If $\mu_u\leq c\mu_v$ on $X$ for a finite $c\geq0$, then $c\geq1$.
\end{corollary}
\begin{proof}
On $\{u<v-\delta\}$ the weighted inequality gives $\mu_u\leq\mathrm{e}^{-\lambda\delta}\mu_v$. Apply strict domination to $v-\delta$ and let $\delta\downarrow0$. If the second assertion failed, strict domination against every translate $v+A$ would contradict finiteness of $u,v$ at a common point.
\end{proof}

\begin{lemma}\label{lem:relative-obstacle-contact}
Let $u,w\in\PSH(X,\theta)$ with $P_\theta(u,w)$ genuine, and put
\[
 P_\theta[u](w)=\left(\lim_{C\to\infty}P_\theta(u+C,w)\right)^*.
\]
Then
\[
 \mu_{P_\theta[u](w)}
 \leq\mathbf1_{\{P_\theta[u](w)=w\}}\mu_w.
\]
\end{lemma}
\begin{proof}
Set $p_C=P_\theta(u+C,w)$ for $C\geq0$ and $p=P_\theta[u](w)$. The envelopes increase, are bounded below by $p_0$ and above by $w$, and converge in capacity by \cref{lem:monotone-psh-capacity}. Assigning ties in the contact set to $w$ gives
\[
 \mu_{p_C}\leq\mathbf1_{\{p_C=u+C<w\}}\mu_u
                    +\mathbf1_{\{p_C=w\}}\mu_w.
\]
The first measure has total mass at most
\[
 \varepsilon_C=\mu_u(\{u\leq\sup_Xw-C\})\longrightarrow0,
\]
because $\mu_u$ is finite and non-pluripolar. Testing against nonnegative continuous functions and applying \cref{lem:weighted-extended-lsc} gives $\mu_p\leq\mu_w$.

The bounded quasi-continuous weight $\chi=\min(1,(w-p)_+)$, defined arbitrarily on the polar indeterminacy set, vanishes on $\{p_C=w\}$, since $p_C\leq p\leq w$. Thus $\int_X\chi\,\mu_{p_C}\leq\varepsilon_C$. Weighted lower semicontinuity gives $\int_X\chi\,\mu_p=0$, so $\mu_p$ is carried by $\{p=w\}$. This proves the asserted restriction of $\mu_w$.
\end{proof}

\begin{theorem}\label{thm:ordinary-envelope-volumes}
For $u\in\PSH(X,\theta)$ and $p=P_\theta[u]$,
\[
 \underline{\operatorname{vol}}_\theta(u)
       =\underline{\operatorname{vol}}_\theta(p).
\]
If this volume is positive, then also
\[
 \overline{\operatorname{vol}}_\theta(u)
       =\overline{\operatorname{vol}}_\theta(p),\qquad
 u\in\mathcal E_{\mathrm{env}}(X,\theta,p),\qquad
 P_\theta[p]=p.
\]
\end{theorem}
\begin{proof}
If $\underline{\operatorname{vol}}_\theta(p)=0$, lower-volume equality follows from monotonicity. Suppose henceforth that it is positive. Normalize $u\leq0$ and choose any $w\simeq p$ in $\PSH(X,\theta)$. Fix $K\geq0$ with $p-K\leq w\leq p+K$. Since $u\leq p$, the rooftop $p_C=P_\theta(u+C,w)$ satisfies $u-K\leq p_C\leq u+C$ for each $C\geq0$. In particular it is genuine and equivalent to $u$. Its increasing canonical limit $r=P_\theta[u](w)$ satisfies
\[
 P_\theta(u+C+K,0)-K\leq p_C\leq w,
 \qquad p-K\leq r\leq w\leq p+K.
\]
Indeed, the leftmost potential lies below $u+C$ and below $p-K\leq w$, so is a rooftop competitor. Thus $r\simeq p$, and \cref{prop:relative-elementary} puts $r$ in $\mathcal E_{\mathrm{env}}(X,\theta,p)$.

By \cref{lem:relative-obstacle-contact}, $\mu_r$ vanishes on $\{r<w\}$. \Cref{thm:big-relative-domination}, with coefficient zero and reference $p$, gives $r\geq w$; hence $r=w$. Capacity convergence and lower semicontinuity now yield
\[
 \int_X\mu_w\leq\liminf_{C\to\infty}\int_X\mu_{p_C}
       \leq\overline{\operatorname{vol}}_\theta(u).
\]
Taking the supremum over $w\simeq p$ proves the upper-volume inequality in the required direction; the reverse inequality is monotonicity. \Cref{thm:relative-equivalence} then gives both relative-class membership and lower-volume equality. Finally \cref{prop:relative-envelope-equality} gives $P_\theta[p]=P_\theta[u]=p$.
\end{proof}

\begin{corollary}\label{cor:full-mass-envelope-characterization}
If $\varphi=P_\theta[\varphi]$ has positive lower volume, then
\[
 \mathcal E_{\mathrm{env}}(X,\theta,\varphi)
 =\mathcal E_{\mathrm{vol}}(X,\theta,\varphi)
 =\{u\in\PSH(X,\theta):P_\theta[u]=\varphi\}.
\]
\end{corollary}
\begin{proof}
The first equality is \cref{thm:relative-equivalence}. Membership implies envelope equality by \cref{prop:relative-envelope-equality}; the converse follows from \cref{thm:ordinary-envelope-volumes}.
\end{proof}

\section{Decreasing model potentials}

The upper volumes of decreasing ordinary model potentials form a Cauchy sequence. The deficit lemma gives convex comparisons with coefficients tending to one; ordinary envelopes normalize the remaining potentials so that compactness passes these comparisons to the limit.

\begin{theorem}\label{thm:decreasing-model-volume-continuity}
Suppose $q_j=P_\theta[q_j]$ decreases to $q\in\PSH(X,\theta)$ and $\underline{\operatorname{vol}}_\theta(q)>0$. Then $P_\theta[q]=q$ and
\[
 \underline{\operatorname{vol}}_\theta(q_j)\searrow
       \underline{\operatorname{vol}}_\theta(q),\qquad
 \overline{\operatorname{vol}}_\theta(q_j)\searrow
       \overline{\operatorname{vol}}_\theta(q).
\]
For every fixed smooth compact modification $\pi:Z\to X$,
\[
 \sup_{z\in Z}\bigl|\nu(q_j\circ\pi,z)-\nu(q\circ\pi,z)\bigr|
       \longrightarrow0.
\]
In particular all generic divisorial Lelong numbers converge.
\end{theorem}
\begin{proof}
Since $q\leq0$, monotonicity of ordinary envelopes gives $q\leq P_\theta[q]\leq P_\theta[q_j]=q_j$. Letting $j\to\infty$ proves the fixed-point assertion. Put
\[
 \delta=\underline{\operatorname{vol}}_\theta(q)>0,
 \qquad L=\lim_j\overline{\operatorname{vol}}_\theta(q_j),
 \qquad e_k=\overline{\operatorname{vol}}_\theta(q_k)-L.
\]
The limit exists by \cref{thm:volume-monotonicity} and bounded mass. After discarding finitely many indices, choose $b_k>1$ tending to infinity with $b_k^n e_k<\delta$, and set $a_k=1-b_k^{-1}$.

\Needspace{7\baselineskip}
For every $j\geq k$, the volume deficit between $q_j$ and $q_k$ is at most $e_k$, whereas the lower volume of $q_k$ is at least $\delta$. Thus \cref{lem:volume-deficit-stretch} supplies $h_{k,j}\in\PSH(X,\theta)$ with
\[
 a_kq_k+(1-a_k)h_{k,j}\leq q_j.
\]
Combining normalized ordinary-envelope competitors, as in the proof of \cref{prop:relative-envelope-equality}, gives
\[
 a_kq_k+(1-a_k)P_\theta[h_{k,j}]\leq q_j.
\]
The potentials $P_\theta[h_{k,j}]$ have supremum zero. For fixed $k$, normalized compactness gives a subsequential $L^1$ limit $h_k$ as $j\to\infty$. Passing to the limit and applying \cref{lem:qpsh-almost-everywhere-order} gives
\[
 a_kq_k+(1-a_k)h_k\leq q.
\]
For either volume, homogeneity, addition and monotonicity now yield
\[
 a_k^n\operatorname{vol}_\theta(q_k)
 \leq\operatorname{vol}_\theta\bigl(a_kq_k+(1-a_k)h_k\bigr)
 \leq\operatorname{vol}_\theta(q)
 \leq\operatorname{vol}_\theta(q_k).
\]
Letting $k\to\infty$ proves both volume limits. Apply \cref{prop:volume-deficit-lelong} with $u=q$, $v=q_j$: its deficit tends to zero and its denominator is at least $\delta$. This proves the uniform Lelong-number convergence on each fixed model.
\end{proof}

We apply this theorem on the submanifold $Y$, returning to the ambient data of \cref{chap:trace-data}. The majorant in a smaller background gives the trace envelope positive lower volume.

\begin{corollary}\label{cor:mass-bound-exact-trace}
Suppose $Y$ has bounded mass. Choose a Hermitian form $\Omega_0\geq\theta|_Y$ and a Hermitian form $\Omega\geq2\Omega_0$ with $v_-(\Omega)>0$. Then
\[
 \nu\bigl(\mathcal T_\Omega(u;Y)\circ\pi,E\bigr)=\tau_{Y,u}(E)
\]
for every prime divisor $E$ on every smooth compact modification $\pi:Z\to Y$.
\end{corollary}
\begin{proof}
Put $T=\mathcal T_\Omega(u;Y)$ and $r=\dim Y$. The normalized majorant $v$ of \cref{prop:trace-profile-majorant} lies below $T$ and belongs to $\PSH(Y,\Omega/2)$. For every $w\simeq T$ in $\PSH(Y,\Omega)$, a constant translate of $v$ lies below $w$. By \cref{lem:smaller-background-mass-lower},
\[
 \int_Y\mathrm{MA}_\Omega(w)\geq2^{-r}v_-(\Omega),
 \qquad
 \underline{\operatorname{vol}}_\Omega(T)
       \geq2^{-r}v_-(\Omega)>0.
\]
The restricted approximants belong to $\PSH(Y,\Omega)$ for all sufficiently large indices. By \cref{thm:restricted-analytic-envelope-limit}, their dyadic ordinary envelopes $q_j=P_\Omega[u_{2^j}|_Y]$ decrease to $T$. The preceding theorem applies on $Y$ and gives
\[
 \nu(T\circ\pi,E)
 =\lim_j\nu(q_j\circ\pi,E)
 =\lim_j2^{-j}\operatorname{ord}_E\mathfrak a_{2^j}(u;Y)
 =\tau_{Y,u}(E),
\]
as required.
\end{proof}

\section{Positive volume and smaller backgrounds}

In a semipositive background, stretching $u$ by a factor $b>1$ and then dividing a competitor by $b$ produces a potential below $u$ in background $b^{-1}\theta$. Positive lower volume allows a small stretch: bounded mass controls its mixed-product error, while contact support excludes collapse of the approximating envelopes.

\Needspace{8\baselineskip}
\begin{proposition}\label{prop:positive-volume-smaller-background}
Assume also that $\theta\geq0$. If $u\in\PSH(X,\theta)$ and
$\underline{\operatorname{vol}}_\theta(u)>0$, then
$P_{\varepsilon\theta}(u)$ is genuine for some $\varepsilon\in(0,1)$.
\end{proposition}
\begin{proof}
Translate $u$ so that $u\leq0$, and put
\[
 \delta=\underline{\operatorname{vol}}_\theta(u),\qquad
 M=\sup_{v\in\PSH(X,\theta)}\int_X\mathrm{MA}_\theta(v)<+\infty.
\]
\Needspace{9\baselineskip}
Choose $b\in(1,2]$ with $2^n(b-1)M<\delta$, and write $s=b-1$.
For integers $j\geq1$, set
\[
 r_j=\max(u,-j/s),\qquad
 f_j=u+s r_j=\max(bu,u-j),\qquad p_j=P_\theta(f_j).
\]
Semipositivity makes $r_j$ a $\theta$-plurisubharmonic function. Since
$u-j\leq f_j\leq u$, the envelopes are genuine, satisfy
$u-j\leq p_j\leq u$ and are equivalent to $u$.

\Needspace{11\baselineskip}
Suppose $P_\theta(bu)\equiv-\infty$. Then $p_j\downarrow-\infty$, and
\cref{thm:normalized-compactness} gives $\sup_Xp_j\to-\infty$.
The contact theorem carries $\mathrm{MA}_\theta(p_j)$ on
$D_j=\{p_j=f_j\}$. Since $f_j$ is $b\theta$-plurisubharmonic,
\cref{prop:ordered-contact-general} and $\theta\leq b\theta$ give
\[
 \mathrm{MA}_\theta(p_j)
 \leq\mathbf1_{D_j}\mathrm{MA}_{b\theta}(f_j).
\]
For $0<s\leq1$, expansion of the positive mixed products gives the measure inequality
\[
 0\leq\mathrm{MA}_{b\theta}(f_j)-\mathrm{MA}_\theta(u)
 \leq 2^n s\,\mathrm{MA}_\theta\bigl((u+r_j)/2\bigr).
\]
Indeed, each mixed term containing $k\geq1$ copies of
$\theta+\ddc r_j$ has coefficient $s^k\leq s$; their sum is bounded by
$s$ times the full power of
$(\theta+\ddc u)+(\theta+\ddc r_j)$.

Fix $C>0$. For large $j$, $p_j\leq-C$ on $X$. On $D_j$ we then have
$bu\leq f_j=p_j\leq-C$, so $u\leq-C/b$. Integrating yields
\[
 \delta\leq\int_X\mathrm{MA}_\theta(p_j)
 \leq\int_{\{u\leq-C/b\}}\mathrm{MA}_\theta(u)+2^n sM.
\]
The first term tends to zero as $C\to+\infty$, because the measure is finite and non-pluripolar. This contradicts the choice of $b$. Thus $P_\theta(bu)$ is genuine, and $b^{-1}P_\theta(bu)$ supplies a competitor for $P_{b^{-1}\theta}(u)$.
\end{proof}

For model references, the positive-volume condition is
\begin{equation}\label{eq:li-model-condition}
 P_\theta[\varphi]=\varphi,\qquad
 \underline{\operatorname{vol}}_\theta(\varphi)>0.
\end{equation}
When $\theta\geq0$, the proposition turns this into the smaller-background condition \eqref{eq:als-model-condition}. If also $v_-(\theta)>0$, the converse follows from \cref{lem:smaller-background-mass-lower}: every potential equivalent to $\varphi$ admits a smaller-background competitor with the same $\varepsilon$, and hence
\[
 \underline{\operatorname{vol}}_\theta(\varphi)
 \geq(1-\varepsilon)^n v_-(\theta)>0.
\]
Without this additional lower bound, the smaller-background argument still gives individual positivity, as proved in \cref{sec:relative-envelope-room}.

\Needspace{11\baselineskip}
\section{Volumes of multiplier-ideal envelopes}

Ideal data determine the upper and lower volumes of their associated envelope. For $u\in\PSH(X,\theta)$, set
\[
 \mathcal V^-_{\mathcal I,\theta}(u)
   =\underline{\operatorname{vol}}_\theta(P_\theta[u]_{\mathcal I}),\qquad
 \mathcal V^+_{\mathcal I,\theta}(u)
   =\overline{\operatorname{vol}}_\theta(P_\theta[u]_{\mathcal I}).
\]
They depend only on the multiplier-ideal class of $u$ in the fixed background. By \cref{prop:i-envelope-gauge} and gauge invariance of singularity volumes, they also depend only on the associated $\ddc$-class and the multiplier-ideal data, rather than on its smooth decomposition. Moreover, \cref{prop:p-below-i,thm:volume-monotonicity} give
\[
 \underline{\operatorname{vol}}_\theta(u)
 \leq\mathcal V^-_{\mathcal I,\theta}(u),\qquad
 \overline{\operatorname{vol}}_\theta(u)
 \leq\mathcal V^+_{\mathcal I,\theta}(u).
\]
The ideal envelope is an ordinary fixed point by \cref{cor:ordinary-preserves-i}. Its positive lower volume lets relative full mass detect equality with the ordinary envelope.

\Needspace{14\baselineskip}
\begin{proposition}\label{prop:conditional-i-characterization}
Let $u\in\PSH(X,\theta)$, and put $q=P_\theta[u]_{\mathcal I}$. Suppose $\underline{\operatorname{vol}}_\theta(q)>0$. Then the following are equivalent:
\[
 P_\theta[u]=P_\theta[u]_{\mathcal I},\qquad
 \overline{\operatorname{vol}}_\theta(u)
      =\mathcal V^+_{\mathcal I,\theta}(u).
\]
If they hold, the corresponding equality for the lower volume also holds.
\end{proposition}
\begin{proof}
If the envelopes agree, \cref{thm:ordinary-envelope-volumes} gives equality of both volumes with those of $q$.

Conversely, upper-volume equality and $u\preceq q$ put $u$ in $\mathcal E_{\mathrm{vol}}(X,\theta,q)$. Since $q=P_\theta[q]$, \cref{cor:full-mass-envelope-characterization} gives $P_\theta[u]=q$.
\end{proof}

Under the hypotheses of the proposition, \cref{cor:dyadic-goodness-criterion} also identifies equality of the upper volumes with convergence of the dyadic ordinary envelopes to $P_\theta[u]$.

\begin{problem}\label{qu:hermitian-i-characterization}
For a big nonclosed background on a manifold with bounded mass, find criteria for $\underline{\operatorname{vol}}_\theta(P_\theta[u]_{\mathcal I})>0$ in terms of restricted multiplier ideals or divisorial orders. When this lower volume vanishes, which implications between envelope equality and upper-volume equality in \cref{prop:conditional-i-characterization} remain valid?
\end{problem}
